\documentclass[a4paper,11pt]{ltjsarticle}

\usepackage[margin=23mm]{geometry}
\usepackage{amsmath,amssymb,amsthm,mathtools}
\usepackage{booktabs}
\usepackage{xcolor}
\usepackage[hidelinks]{hyperref}
\usepackage{enumitem}

\definecolor{ink}{HTML}{1F2933}
\definecolor{accent}{HTML}{2F5D62}
\definecolor{soft}{HTML}{EEF6F4}

\hypersetup{
  colorlinks=true,
  linkcolor=accent,
  urlcolor=accent,
  citecolor=accent
}

\setlist[itemize]{leftmargin=1.5em,itemsep=0.25em,topsep=0.25em}
\setlength{\parindent}{1em}
\setlength{\parskip}{0.2em}

\theoremstyle{definition}
\newtheorem{definition}{定義}[section]
\newtheorem{example}{例}[section]

\theoremstyle{plain}
\newtheorem{proposition}{命題}[section]
\newtheorem{theorem}{定理}[section]

\theoremstyle{remark}
\newtheorem*{remark}{注意}

\newcommand{\HN}{H_N}
\newcommand{\V}{\mathrm{V}}
\newcommand{\E}{\mathrm{E}}
\newcommand{\tree}{\tau}
\newcommand{\dd}{\mathrm{d}}

\title{\vspace{-1.2cm}
{\Large\color{accent}Hanoi graph 入門}\\[0.4em]
{\huge サイズ・状態数・最短解長・全域木エントロピー}\\[0.4em]
{\large 明日の発表のための短い教科書}}
\author{}
\date{}

\begin{document}
\color{ink}
\maketitle
\vspace{-1.0cm}

\begin{center}
\begin{minipage}{0.86\linewidth}
\small
本ノートは，Hanoi graph のうち発表で扱う最初の部分だけをまとめる。
中心に置く量は，円盤数 \(N\)，状態数 \(3^N\)，標準問題の最短解長
\(2^N-1\)，および Zhang--Wu--Li--Comellas による全域木数と
全域木エントロピー
\[
  h=\frac{1}{4}\log 15
\]
である。記述は Wikipedia の Hanoi graph 項目と
Zhang et al. の論文に沿う。
\end{minipage}
\end{center}

\section{Hanoi graph とは何か}

Tower of Hanoi の各配置をひとつの点とみなし，合法手で移れる配置どうしを線で結ぶ。
このようにして得られる無向グラフが Hanoi graph である。

\begin{definition}[Hanoi graph]
\(k\) 本の杭と \(N\) 枚の円盤からなる Tower of Hanoi を考える。
各頂点は，それぞれの円盤がどの杭にあるかを表す状態である。
ふたつの頂点を結ぶ辺は，一手の合法移動で互いに移れることを表す。
このグラフを
\[
  H_k^N
\]
と書く。特に，古典的な3本杭の場合を
\[
  \HN := H_3^N
\]
と書く。
\end{definition}

ここで重要なのは，状態を「円盤の並び」としてではなく，
「各円盤がどの杭にいるか」という情報として記述できることである。
円盤の大小関係により，各杭の上では自動的に小さい円盤が上，大きい円盤が下になる。
したがって，各円盤について選択肢は杭の数だけでよい。

\begin{proposition}[状態数]
\(k\) 本杭・\(N\) 円盤の Hanoi graph は
\[
  |V(H_k^N)|=k^N
\]
個の頂点を持つ。特に3本杭では
\[
  |V(\HN)|=3^N.
\]
\end{proposition}

\begin{proof}
各円盤は \(k\) 本の杭のどれかに置かれている。
円盤は \(N\) 枚あるので，状態は長さ \(N\) の列
\[
  (a_1,a_2,\ldots,a_N), \qquad a_i\in\{1,\ldots,k\}
\]
として表せる。よって状態数は \(k^N\) である。
\end{proof}

\begin{example}[小さい \(N\) の状態数]
3本杭の場合，
\[
  N=1:\ 3,\qquad
  N=2:\ 9,\qquad
  N=3:\ 27,\qquad
  N=4:\ 81.
\]
円盤が1枚増えるごとに状態数は3倍になる。
\end{example}

\section{辺と自己相似構造}

Hanoi graph は単に状態を並べたグラフではなく，きれいな再帰構造を持つ。
最大の円盤がどの杭にあるかで状態を分類すると，
\(\HN\) は \(H_{N-1}\) の3個のコピーから作られる。

\begin{proposition}[3本杭の場合の辺数]
3本杭 Hanoi graph \(\HN\) の辺数は
\[
  |E(\HN)|=\frac{3(3^N-1)}{2}
\]
である。
\end{proposition}

\begin{proof}
Wikipedia の Hanoi graph 項目では，一般の \(k\) 本杭について
\[
  |E(H_k^N)|
  =
  \frac{1}{2}\binom{k}{2}\left(k^N-(k-2)^N\right)
\]
と与えられている。ここで \(k=3\) とおくと，
\[
  |E(\HN)|
  =
  \frac{1}{2}\binom{3}{2}(3^N-1^N)
  =
  \frac{3(3^N-1)}{2}.
\]
\end{proof}

\begin{remark}
3本杭では，多くの内部状態は次数3を持つ。一方，すべての円盤が同じ杭に積まれている
3つの角の状態では，動かせる手は2通りである。グラフ全体としては，
三角形を再帰的に3分割していく Sierpiński triangle に似た形を持つ。
\end{remark}

\section{標準問題の最短解長}

標準的な Tower of Hanoi では，すべての円盤がある杭に積まれた状態から，
別の杭へすべて移す。Hanoi graph で見ると，これは三角形の角から別の角への
最短路を求める問題である。

\begin{theorem}[標準解長]
3本杭・\(N\) 円盤の標準 Tower of Hanoi の最短手数は
\[
  L^*(N)=2^N-1
\]
である。Wikipedia の Hanoi graph 項目では，3本杭 Hanoi graph の直径も
この値であると述べられている。
\end{theorem}

\begin{proof}[再帰による説明]
最大円盤を出発杭から目的杭へ動かすには，その上にある \(N-1\) 枚をいったん
補助杭へ移しておく必要がある。そのために少なくとも \(L^*(N-1)\) 手が必要である。
次に最大円盤を1手で目的杭へ移す。最後に，補助杭にある \(N-1\) 枚を目的杭へ
移すために再び \(L^*(N-1)\) 手が必要である。したがって
\[
  L^*(N)=2L^*(N-1)+1,\qquad L^*(1)=1.
\]
この漸化式を解くと
\[
  L^*(N)=2^N-1
\]
を得る。
\end{proof}

ここで，状態数 \(3^N\) と最短解長 \(2^N-1\) は異なる指数増大である。
状態空間は3倍ずつ広がるが，標準問題の最短路はおよそ2倍ずつ長くなる。
この違いは，発表で強調すると理解しやすい。

\begin{center}
\begin{tabular}{rrrr}
\toprule
\(N\) & \(|V(\HN)|=3^N\) & \(L^*(N)=2^N-1\) & \(|E(\HN)|\) \\
\midrule
1 & 3 & 1 & 3 \\
2 & 9 & 3 & 12 \\
3 & 27 & 7 & 39 \\
4 & 81 & 15 & 120 \\
5 & 243 & 31 & 363 \\
\bottomrule
\end{tabular}
\end{center}

\section{全域木: グラフ全体の接続構造を数える}

次に Zhang, Wu, Li, Comellas の論文
\textit{The number and degree distribution of spanning trees in the Tower of Hanoi graph}
の主結果を見る。ここで扱うのは，解法の数ではなく，Hanoi graph の上に張れる
全域木の数である。

\begin{definition}[全域木]
有限連結グラフ \(G=(V,E)\) の全域木とは，すべての頂点を含み，かつ閉路を持たない
連結部分グラフである。全域木数を \(\tau(G)\) と書く。
\end{definition}

全域木は，グラフ全体を「枝だけで，しかし全頂点をつないだ」構造である。
Hanoi graph の全域木数は，状態空間そのものの大きさとは別の，
接続構造の多様性を表す。

\begin{theorem}[Zhang et al. の全域木数公式]
\(N\) 円盤の3本杭 Hanoi graph \(\HN\) の全域木数を
\[
  s_N := \tau(\HN)
\]
とする。このとき
\[
  s_N
  =
  3^{\frac{1}{4}3^N+\frac{1}{2}N-\frac{1}{4}}
  5^{\frac{1}{4}3^N-\frac{1}{2}N-\frac{1}{4}}.
\]
\end{theorem}

この式は，Hanoi graph の再帰構造から得られる。
論文では，単に全域木だけを直接数えるのではなく，いくつかの種類の全域部分グラフの
個数の間に漸化式を立て，それを解くことで閉じた形の公式を導いている。

\begin{example}[小さい \(N\) での確認]
公式から，
\[
\begin{aligned}
s_1
&=
3^{\frac{3}{4}+\frac{1}{2}-\frac{1}{4}}
5^{\frac{3}{4}-\frac{1}{2}-\frac{1}{4}}
=3,\\
s_2
&=
3^{\frac{9}{4}+1-\frac{1}{4}}
5^{\frac{9}{4}-1-\frac{1}{4}}
=3^3\cdot 5=135.
\end{aligned}
\]
\(H_1\) は三角形なので，全域木は3本の辺から1本を除く3通りであり，
\(s_1=3\) と一致する。
\end{example}

\section{全域木エントロピー}

全域木数 \(s_N\) は非常に速く増える。
その増え方を頂点数あたりで正規化したものが，全域木エントロピーである。

\begin{definition}[全域木エントロピー]
\[
  h
  =
  \lim_{N\to\infty}
  \frac{\log s_N}{|V(\HN)|}
  =
  \lim_{N\to\infty}
  \frac{\log s_N}{3^N}.
\]
\end{definition}

\begin{theorem}[Hanoi graph の全域木エントロピー]
Zhang et al. の公式から，
\[
  h
  =
  \frac{1}{4}\log 15
  \approx 0.677.
\]
\end{theorem}

\begin{proof}
全域木数公式の対数を取ると，
\[
\begin{aligned}
\log s_N
&=
\left(\frac{1}{4}3^N+\frac{1}{2}N-\frac{1}{4}\right)\log 3\\
&\quad+
\left(\frac{1}{4}3^N-\frac{1}{2}N-\frac{1}{4}\right)\log 5.
\end{aligned}
\]
これを \(3^N\) で割ると，
\[
\frac{\log s_N}{3^N}
=
\frac{1}{4}(\log 3+\log 5)
+ \frac{N}{2\cdot 3^N}(\log 3-\log 5)
- \frac{1}{4\cdot 3^N}(\log 3+\log 5).
\]
\(N/3^N\to 0\) かつ \(1/3^N\to0\) なので，
\[
  h
  =
  \frac{1}{4}(\log 3+\log 5)
  =
  \frac{1}{4}\log 15.
\]
\end{proof}

この結果は，次のように読むとよい。
\[
  \log s_N \sim h\,3^N,
  \qquad
  s_N \sim \exp(h\,3^N).
\]
つまり，頂点数 \(3^N\) を基準にすれば，全域木数は通常の指数増大である。
一方で，円盤数 \(N\) を基準に見ると，指数の中にさらに \(3^N\) が現れる。
このため，Hanoi graph の接続構造の多様性は，円盤数に対して非常に急激に増える。

\section{発表での言い方}

\begin{itemize}
\item Hanoi graph \(\HN\) は，3本杭・\(N\) 円盤 Tower of Hanoi の全状態を頂点にしたグラフである。
\item 各円盤が3本の杭のどこにあるかを選ぶので，状態数は
\[
  |V(\HN)|=3^N.
\]
\item 標準問題の最短解長は
\[
  L^*(N)=2^N-1.
\]
状態空間の増え方 \(3^N\) と，最短解長の増え方 \(2^N\) は別物である。
\item Zhang et al. は \(\HN\) の全域木数を厳密に求めた。
\[
  s_N
  =
  3^{\frac{1}{4}3^N+\frac{1}{2}N-\frac{1}{4}}
  5^{\frac{1}{4}3^N-\frac{1}{2}N-\frac{1}{4}}.
\]
\item その結果，全域木エントロピーは
\[
  h=\lim_{N\to\infty}\frac{\log s_N}{3^N}
  =\frac{1}{4}\log 15
  \approx0.677.
\]
\end{itemize}

\section*{参考文献}

\begin{enumerate}[leftmargin=1.7em]
\item Wikipedia, ``Hanoi graph.'' \\
\url{https://en.wikipedia.org/wiki/Hanoi_graph}
\item Zhongzhi Zhang, Shunqi Wu, Mingyun Li, Francesc Comellas,
``The number and degree distribution of spanning trees in the Tower of Hanoi graph,''
\textit{Theoretical Computer Science} 609, 443--455, 2016. arXiv:1510.07949.\\
\url{https://arxiv.org/abs/1510.07949}
\end{enumerate}

\end{document}
